\section{Hilbert--valued folded Parseval and tomography spectrum}
\label{sec:tpc39-folded}

The minimal quotient deliberately identifies some positive and negative
aliases.  This section computes that folding exactly.

\subsection{General quotient sizes}

Let
\begin{equation}
 L<r\le2L+1
 \label{eq:tpc39-general-r-range}
\end{equation}
and put \(\zeta_r=\ee{1/r}\).  For Hilbert columns
\(Z_{-L},\ldots,Z_L\in\cH\), define
\begin{align}
 Y_\nu^{(r)}
 &:=\sum_{k=-L}^{L}\zeta_r^{\nu k}Z_k,
 \qquad 0\le\nu<r,
 \label{eq:tpc39-general-twists}\\
 W_a^{(r)}
 &:=\sum_{\substack{-L\le k\le L\\k\equiv a\pmod r}}Z_k,
 \qquad a\in\Z/r\Z.
 \label{eq:tpc39-residue-buckets}
\end{align}

\begin{theorem}[Hilbert--valued folded Parseval]
\label{thm:tpc39-folded-Parseval}
For every \(r\) in \eqref{eq:tpc39-general-r-range},
\begin{equation}
 \boxed{
 \frac1r\sum_{\nu=0}^{r-1}\|Y_\nu^{(r)}\|^2
 =\sum_{a\bmod r}\|W_a^{(r)}\|^2.}
 \label{eq:tpc39-folded-Parseval}
\end{equation}
At the minimal value \(r=q=L+1\),
\begin{equation}
 \boxed{
 \frac1q\sum_{\nu=0}^{q-1}\|Y_\nu^{(q)}\|^2
 =\|Z_0\|^2+
 \sum_{a=1}^{L}\|Z_a+Z_{a-q}\|^2.}
 \label{eq:tpc39-minimal-folded-Parseval}
\end{equation}
\end{theorem}

\begin{proof}
Expand the squared norms and average the phase
\(\zeta_r^{\nu(k-k')}\).  Additive orthogonality retains precisely the
pairs \(k\equiv k'\pmod r\), proving
\eqref{eq:tpc39-folded-Parseval}.  If \(r=L+1\), the zero residue class is
\(\{0\}\), while residue \(a\), \(1\le a\le L\), contains the pair
\(\{a,a-q\}\).
\end{proof}

\subsection{Complete singular spectrum}

Define the normalized tomography operator
\begin{equation}
 \cT_r:\cH^{2L+1}\longrightarrow\cH^r,
 \qquad
 (\cT_rZ)_\nu
 =r^{-1/2}\sum_{k=-L}^{L}\zeta_r^{\nu k}Z_k.
 \label{eq:tpc39-tomography-operator}
\end{equation}
Its scalar Gram matrix is
\begin{equation}
 (\cT_r^*\cT_r)_{k,k'}
 =\one_{k\equiv k'\pmod r}.
 \label{eq:tpc39-tomography-Gram}
\end{equation}

\begin{theorem}[Alias--index tomography spectrum]
\label{thm:tpc39-tomography-spectrum}
For \(L<r\le2L+1\), the spectrum of the scalar alias--index Gram matrix
in \eqref{eq:tpc39-tomography-Gram} is
\begin{equation}
 \boxed{
 \begin{array}{c|c}
 \text{eigenvalue}&\text{multiplicity}\\ \hline
 2&2L+1-r\\
 1&2r-(2L+1)\\
 0&2L+1-r.
 \end{array}}
 \label{eq:tpc39-general-spectrum}
\end{equation}
At \(r=q=L+1\), this becomes
\begin{equation}
 2\ (L\text{ times}),
 \qquad1\ (1\text{ time}),
 \qquad0\ (L\text{ times}).
 \label{eq:tpc39-minimal-spectrum}
\end{equation}
At the minimal value \(r=q\), the eigenvalue--one subspace is the isolated
zero column.  The kernel is
\begin{equation}
 Z_0=0,
 \qquad
 Z_a=-Z_{a-q},
 \quad1\le a\le L.
 \label{eq:tpc39-minimal-kernel}
\end{equation}
As an operator on \(\cH^{2L+1}\), one has
\(\cT_r^*\cT_r=G_r\otimes I_{\cH}\), where \(G_r\) is this scalar
matrix.  Thus for finite \(\dim\cH=d\), every displayed multiplicity is
multiplied by \(d\); the table records the intrinsic alias--index
multiplicities.
\end{theorem}

\begin{proof}
Because \(r>L\), every residue bucket in
\eqref{eq:tpc39-residue-buckets} has size one or two.  There are
\(2L+1-r\) doubletons and \(2r-(2L+1)\) singletons.  On a bucket of size
\(s\), the Gram block in \eqref{eq:tpc39-tomography-Gram} is the
all--ones matrix \(J_s\), whose eigenvalues are \(s\) once and zero
\(s-1\) times.  This proves \eqref{eq:tpc39-general-spectrum}.  The
minimal specialization follows from the explicit buckets in the proof of
\cref{thm:tpc39-folded-Parseval}.
\end{proof}

At the other endpoint \(r=2L+1\), every bucket is a singleton and
\begin{equation}
 \frac1{2L+1}\sum_{\nu=0}^{2L}\|Y_\nu^{(2L+1)}\|^2
 =\sum_{k=-L}^{L}\|Z_k\|^2.
 \label{eq:tpc39-full-DFT-Parseval}
\end{equation}
Thus \(L+1\) twists are sufficient to recover the zero column, whereas
\(2L+1\) twists are required to separate every alias column by this
cyclic construction.

\subsection{Exact position between trace and Gram gates}

Put
\begin{align}
 E_{\rm tr}(Z)
 &:=\sum_{k=-L}^{L}\|Z_k\|^2,
 \label{eq:tpc39-trace-energy}\\
 E_{\rm fold}(Z)
 &:=\|Z_0\|^2+
 \sum_{a=1}^{L}\|Z_a+Z_{a-q}\|^2.
 \label{eq:tpc39-fold-energy}
\end{align}
Then
\begin{equation}
 \|Z_0\|^2\le E_{\rm fold}(Z)\le2E_{\rm tr}(Z).
 \label{eq:tpc39-fold-trace-comparison}
\end{equation}
If \(\mathsf G=(\langle Z_k,Z_{k'}\rangle)\), then
\begin{equation}
 E_{\rm tr}(Z)=\tr\mathsf G
 \le(2L+1)\|\mathsf G\|_{\op}.
 \label{eq:tpc39-Gram-trace-comparison}
\end{equation}
Consequently,
\begin{equation}
 \text{alias Gram}
 \Longrightarrow
 \text{raw trace}
 \Longrightarrow
 \text{minimal folded energy}
 \Longrightarrow
 \text{zero--column energy}.
 \label{eq:tpc39-gate-hierarchy}
\end{equation}

The reverse implications fail in the unrestricted Hilbert class.  Taking
\(Z_a=Tv\), \(Z_{a-q}=-Tv\), and all other columns zero makes
\(E_{\rm fold}=0\) while the trace and Gram norm grow like \(T^2\).
Taking the same pair with equal signs and \(Z_0=0\) makes the zero--column
energy vanish while \(E_{\rm fold}\) is arbitrarily large.  Hence the
minimal tomography bank discards exactly a large collection of directions
that are irrelevant to equality recovery; it is not a disguised full
trace theorem.
